% sections/00_executive_summary.tex  —  Executive Summary (unnumbered)
% \input{} into main.tex; do NOT wrap in \begin{document}.

\section*{Executive Summary}
\addcontentsline{toc}{section}{Executive Summary}

This report presents a \emph{methodological plan} for classifying
single-port $S_{11}$ reflection-coefficient measurements acquired from an
RF sensor operating in the \SIrange{8}{12}{\giga\hertz} band.  The sensor
is exposed to three distinct surface states—\texttt{dry},
\texttt{wet}, and \texttt{ice}—each of which alters the dielectric
environment and therefore produces a characteristic frequency-dependent
signature in the reflected signal.  The ultimate objective is to deploy
the trained classifier on resource-constrained edge hardware, where both
memory and inference latency are critical constraints.

\paragraph{Dataset.}
The available corpus comprises \num{1603} traces distributed across three
class states: \texttt{dry} (\num{601} traces), \texttt{wet} (\num{501} traces),
and \texttt{ice} (\num{501} traces). Each trace records \num{4001} uniformly spaced
frequency points spanning approximately \SIrange{8.0}{12.0}{\giga\hertz}.
Crucially, the dataset uses a \textbf{Cartesian representation} mapping each
point to its \textbf{Real} and \textbf{Imaginary} components
($\text{Re}[S_{11}]$, $\text{Im}[S_{11}]$). This choice explicitly avoids the
``phase wrapping'' discontinuities inherent to polar representations
(Magnitude/Phase), significantly improving numerical stability for
gradient-based learning.

\paragraph{Dual-route representation strategy.}
Two complementary data representations are proposed and systematically
compared throughout this report:
\begin{description}
    \item[Route~A — 1D Spectral:]  The Cartesian frequency-domain vectors are
        retained as one-dimensional feature sequences. The primary encoding is a
        dual-channel $\mathbb{R}^{4001 \times 2}$ vector containing the Real and
        Imaginary components. This route preserves full numerical fidelity and is
        compatible with both classical tabular models and 1D neural architectures.
    \item[Route~B — 2D Image-Like:]  The spectral curves are rasterised
        into fixed-resolution images, enabling the use of established
        convolutional architectures designed for image classification.
        This route trades information fidelity for access to
        transfer-learning ecosystems but introduces rendering artefacts
        that are not present in the measurement physics.
\end{description}

\paragraph{Tiered Candidate Selection.}
Rather than declaring a single architecture outright, this methodology proposes a
\textbf{tiered selection strategy} traversing from minimal-parameter models to
deep architectures. Models are evaluated progressively, deploying the simplest
model that meets accuracy requirements:

\begin{description}
    \item[Level 0 (Unsupervised Probe):] \textbf{K-Means} validates intrinsic
        data separability before supervised training.
    \item[Level 1 (Linear Baselines):] \textbf{Logistic Regression} and
        \textbf{Linear SVM}. These are highly interpretable models whose forward
        pass requires mere dot products. They establish the performance floor.
    \item[Level 2 (Tabular Powerhouses):] \textbf{Gradient Boosted Trees (GBT)},
        \textbf{Random Forest}, and \textbf{RBF SVM}. GBT is a standout champion
        here due to its consistent dominance on tabular data and invariant scaling,
        making it highly edge-viable without floating-point matrix arithmetic.
    \item[Level 3 (Deep Architectures):] \textbf{1D CNN}, \textbf{MLP},
        \textbf{ResNet-1D}, and \textbf{KAN}. The \textbf{1D CNN} is the primary
        champion of this tier due to its perfect inductive alignment with local
        spectral resonances and extremely efficient parameter-sharing footprint.
\end{description}
